\chapter{Discussion}
The evaluation suggests that the overall idea of using an AI assistant to support robotic program comprehension and modification is promising, but the current implementation is not yet a robust solution for practical use. The user study indicates that participants generally understood the value of AI assistance and were willing to try it, but they also encountered substantial friction when the system had to reason over the raw \texttt{.urpx} representation and when the assistant needed to generate concrete editing guidance.

The most positive result is that comprehension correctness remained high in the proposed system. This indicates that the combination of structured program views, retrieval, and assistant explanations can support the basic understanding of robotic programs. However, the lower confidence ratings and the longer completion times show that this support was not frictionless. In other words, the system often helped participants arrive at the correct answer, but it did not make the task feel easier or faster. A large part of the observed time increase appears to come from assistant latency and the interaction overhead of working through chat rather than from the comprehension task itself.

For modification and reuse, the picture is less favorable. The proposed system clearly underperformed the baseline on correctness and completion rate, and the gap was especially visible for tasks that required structural edits. Participants handled parameter updates reasonably well, but tasks that required inserting or restructuring program nodes exposed a mismatch between the system's current interaction model and the actual editing problem. This is consistent with the requirement analysis: the system is expected to support program editing, but the current implementation does so through low-level \texttt{.urpx} editing, which has shown to be a poor fit for many users.

The failure on MT-5 (add wait node) is particularly important. None of the participants completed the wait-node insertion task successfully in the proposed system. This reflects a system limitation. The assistant did not have sufficient knowledge about the exact node structure, and the knowledge base did not contain the required information for this task. In practice, this means the system could not provide the actionable information needed to complete the edit. The result highlights a core dependency of retrieval-augmented systems: if the relevant knowledge is missing, incomplete, or poorly retrieved, the model cannot reliably compensate by generating the right answer on its own.

These findings suggest that the current RAG-style setup works better as a support mechanism for documentation lookup and guided comprehension than as the primary mechanism for code generation and structural modification. For the present domain, the assistant is effectively being used as a one- or few-shot prompt interface over retrieved examples. That is a useful pattern for natural language question answering, but it is a weak foundation for precise program transformation when the target representation is rigid and under-documented. A different editing strategy may therefore be required. Depending on the target user group, a higher-level UI that resembles the teach pendant may be more suitable for novice users, while direct URScript editing may be more appropriate for advanced users. In either case, the edit model should expose domain concepts at a higher level than raw \texttt{.urpx} nodes.

The two conditions, the baseline and the proposed system, differed on multiple dimensions simultaneously, so it is not possible to attribute observed differences in correctness or time to the assistant alone. The baseline condition offered a visual, teach-pendant-like interface, while the proposed system combined JSON editing with chat-based assistance. The observed differences therefore reflect the combined workflow rather than isolating the assistant's contribution. This is an important caveat when interpreting the results: the assistant may have provided useful information, but the overall workflow was more cumbersome than the baseline, which likely contributed to lower confidence and longer completion times. On the other hand, the assistant may have provided misleading information in some cases, which could have contributed to incorrect answers or lower user confidence in their own solutions. The study design does not allow us to disentangle these effects, so the results should be interpreted as reflecting the overall user experience of the proposed system rather than the assistant's performance in isolation.

The results further indicate that response quality depends strongly on the retrieved context. When the retrieval pipeline surfaced relevant material, the assistant was often understandable and occasionally useful. When the retrieved material was sparse or tangential, the assistant produced weaker answers. This makes data coverage and retrieval quality central design concerns. More data, better chunking, and improved similarity search are likely to have a direct effect on usefulness.

At the same time, the current approach raises the question of whether fine-tuning would be a better fit than pure RAG for code-oriented tasks. Fine-tuning a general coding model on a domain-specific URScript or robot-program dataset would require substantial upfront effort, but it could improve structural editing, domain terminology, and output consistency. RAG would still be valuable for domain documentation and project-specific retrieval, so the two approaches are not mutually exclusive. A hybrid design seems most plausible: use retrieval for factual and project-specific context, and use a specialized or fine-tuned model for code generation and modification.

Finally, the response-time results are a practical warning. Even if the outputs are useful, a slow assistant will be difficult to adopt in everyday programming workflows. Some of the latency stems from retrieval and model inference, but part of the problem is also the user experience: participants have to wait for the full answer before they can continue. This suggests that both system optimization and interface changes are needed. Streaming responses, exposing progress indicators, and reducing orchestration overhead would improve perceived responsiveness even before the underlying model stack is fully optimized.

\section{Implications for the Requirements}
The evaluation provides partial support for the requirements defined earlier in the thesis, but also shows where the requirements remain unmet in practice.

For the functional requirements, the strongest support is for requirements related to comprehension and question answering. The study results indicate that FR5 and FR7 are partially satisfied, since the assistant could often answer natural-language questions and participants achieved near-perfect comprehension correctness. FR6 also appears largely satisfied because users generally rated response clarity positively. FR10 is only partially satisfied: in several cases participants perceived low factual correctness or relevance, indicating that uncertainty handling and reliability were not consistently strong.

Requirements related to modification and direct editing support are less well satisfied. FR8 is not sufficiently satisfied in its current form, as modification correctness dropped substantially in the proposed condition. FR4 is also only partially satisfied from a usability perspective: although an editor exists, editing raw \texttt{.urpx} structures introduced substantial friction for structural changes. Requirements tied to overall repository workflow, namely FR1-FR3, were not the primary focus of the controlled tasks, so the current evaluation provides only limited evidence for their practical effectiveness under realistic usage.

FR9 is partially supported. The system retains session context, but qualitative findings indicate that context was not always perceived as relevant by users and may sometimes have introduced irrelevant carry-over. This suggests that conversational memory exists technically, but context selection and grounding need improvement.

The non-functional requirements are also mixed. NFR1 (response within 10 seconds) is not met in the observed user workflow, where assistant response times were substantially higher. NFR2 (modifiability through isolation of LLM functionality) is architecturally supported by the system design. NFR3 (observability) is well-supported through detailed logging, and NFR4 (traceability) is supported as retrieved references are exposed.

For usability-oriented requirements, NFR5 is the most supported because response clarity was generally rated positively. In contrast, NFR6 and NFR7 are not supported by the present results, since perceived usefulness for comprehension and especially modification was limited. NFR8 and NFR9 are also not sufficiently satisfied, as post-study ratings frequently questioned response relevance and factual correctness.

\section{Threats to Validity}
Several limitations affect the strength of the results. The sample size was only five participants, which is too small to support broad generalization. The study also evaluated the system as a whole rather than isolating the assistant from the editor and retrieval pipeline, so the measured differences reflect the combined workflow.

The task set was narrow. The modification tasks focused on relatively basic parameter changes and node insertions, so the results do not reveal how the system would perform on more complex programming tasks. In addition, the tasks were performed in a fixed order within each condition, which may have introduced order effects even though the overall condition order was counterbalanced.

The timing data should also be interpreted with caution because task durations were measured manually with a stopwatch. This introduces measurement noise and makes small time differences difficult to interpret precisely.

The evaluation did not directly score assistant response correctness. A participant could use the assistant, receive an unhelpful answer, and still complete the task independently. In that case, the study records both assistant use and task success, but it does not distinguish how much the assistant actually contributed to the outcome.


\section{Future Work} \label{sec:discussion:future-work}
Future work should focus on improving the assistant along four dimensions: data coverage, retrieval quality, interaction design, and evaluation methodology.

\subsection{Data and Knowledge Base}
The most direct improvement is to expand the knowledge base. More robot programs, more documentation, and more examples of domain-specific structures would make retrieval more reliable. In addition to official Universal Robots documentation, useful sources would include partner material, customer examples, support forums, and community discussions such as Discord or other Q\&A channels. The assistant would likely benefit from more question-answer pairs, more annotated program examples, and better coverage of uncommon node types and editing patterns.

\subsection{Retrieval and Prompting}
Better retrieval can improve response quality. Better chunking strategies, hybrid similarity search, and improved query rewriting may all improve the relevance of retrieved context. Prompt engineering is another low-cost opportunity: better system prompts, clearer task framing, and more explicit instructions for uncertainty handling could improve response quality without changing the overall architecture.

In addition, the orchestration logic should be improved, so the assistant makes better decisions about when to retrieve, which source type to query, and how many tool steps are actually needed before synthesis. The current behavior suggests that unnecessary or weakly targeted tool calls can increase latency and sometimes reduce response quality by introducing noisy context.

Memory management is also an important improvement area. Future versions should use stronger context-selection policies, so prior conversation history does not over-pollute new responses with stale or irrelevant details.

At the same time, a pure RAG setup may not be the best long-term solution for code generation and program modification. A hybrid approach may be more promising: use retrieval for documentation and project-specific facts, and use a specialized model for code-oriented reasoning. A fine-tuned model trained on URScript or a similar domain-specific corpus could be more reliable for edits, especially when the output must follow strict structural constraints. Such a model would require a suitable dataset and more training effort, but it could reduce the dependence on fragile few-shot prompting.

\subsection{Editing Workflow and User Interface}
The editing workflow should be reconsidered. Raw \texttt{.urpx} editing is a poor interaction model for many users because the representation is too low-level for novice users and too brittle for structural edits. Future versions should either provide a higher-level interface inspired by the teach pendant or support direct URScript editing for advanced users. Which option is preferable depends on the intended user group, but both would likely be more usable than editing the raw program file.

More agentic behavior would elevate the assistant's capabilities: the assistant should be able to inspect the current program, propose edits, validate them, and ideally help execute or simulate the programs.

An integration with a language server protocol could also be useful if the editing workflow is based on text or structured program files. A language server could provide syntax awareness, hover information, references, diagnostics, and code actions. That would make the assistant less dependent on ad hoc prompting and would ground its suggestions in the structure of the actual program format.

\subsection{Performance and Responsiveness}
Improving response time is important for usability. One option is to use a smaller and more specialized orchestration model, reserving larger models only for tasks that actually need them. Another option is to stream more intermediate progress and the final answer, so users do not have to wait while retrieval and generation are running. If the system displays intermediate progress, users may better tolerate latency even when the underlying response time remains nontrivial.

Another direction of improvement is to move away from a local-only model setup. While local deployment is useful for early prototyping, future versions should evaluate hybrid or cloud-hosted model serving for demanding tasks. This could provide access to stronger models, faster inference, and more stable response quality. A practical approach is a hybrid routing strategy in which routine requests are handled locally, while complex requests are escalated to hosted models when higher capability is needed.

\subsection{Evaluation}
Future evaluations should use a larger and more diverse participant sample to improve generalizability. Task timing should be automated rather than measured manually. The task order should also be counterbalanced within each condition so that task-specific order effects are reduced.

The evaluation should additionally measure assistant response correctness directly, not only final task correctness. This would make it possible to separate cases where the assistant was helpful from cases where the participant succeeded despite the assistant. Future studies should also include a broader task set, including more complex edits and more realistic programming scenarios, so that the limits of the system are better understood.
